\subsubsection{Deriving the Laughlin/Jastrow form}
The following reduction makes the boson construction a directly
reproducible example of the conformal-block method. Let
$O=\{i_1<\cdots<i_M\}$ be the occupied sites and write $d_{ij}=z_i-z_j$
for $i<j$.
\begin{enumerate}
\item Expand the site-charge exponent:
\begin{equation}
 \frac{s_i s_j}{2}=2n_i n_j-n_i-n_j+\frac12.
\end{equation}
The last term contributes the same factor to every configuration and
can be absorbed into normalization. The first term produces
$\prod_{\mu<\nu}(z_{i_\mu}-z_{i_\nu})^2$.
\item Collect the two terms linear in occupation into one factor per
occupied site. Define
\begin{equation}
 D_i=\prod_{j<i}(z_j-z_i)\prod_{j>i}(z_i-z_j).
\end{equation}
Then the configuration-dependent pair product is
$\prod_{\mu<\nu}(z_{i_\mu}-z_{i_\nu})^2\prod_{i\in O}D_i^{-1}$.
The endpoint factors similarly reduce to
$\prod_{i\in O}(-z_i)^a(1-z_i/W)^b$.
\item Use the fact that the $z_j$ are the roots of $z^N-1$:
\begin{equation}
 \prod_{j\ne i}(z_i-z_j)=Nz_i^{N-1},\qquad
 D_i=(-1)^iNz_i^{N-1}.
\end{equation}
The factors $(-1)^i$ cancel the Marshall signs. Since $z_i^N=1$,
$z_i^{-(N-1)}=z_i$. For integer $a$, the remaining
$(-1)^{aM}$ is a configuration-independent phase.
\item The explicit polynomial amplitude is therefore
\begin{align}
 \psi_{a,b}(O)&=\frac{\mathcal P_{a,b}(O;W)}{\sqrt{\mathcal Z_{a,b}(W)}},
 \notag\\
 \mathcal P_{a,b}(O;W)&=
 \prod_{\mu<\nu}(z_{i_\mu}-z_{i_\nu})^2
 \prod_{\mu=1}^{M}z_{i_\mu}^{a+1}(1-z_{i_\mu}/W)^b ,
 \label{eq:bosonpolynomial}
\end{align}
where $\mathcal Z_{a,b}(W)=\sum_{|O|=M}|\mathcal P_{a,b}(O;W)|^2$.
An overall phase is immaterial. Equation~\eqref{eq:bosonpolynomial}
is the exponent-two Laughlin/Jastrow form used by the original method.
At strict infinity, set every $(1-z_i/W)^b$ to one.
\end{enumerate}
Although $b$ disappears from the explicit position factors at infinity,
it still fixes $M=(N-a-b)/2$. Thus changing $b$ changes the state.
The block and polynomial formulas give the same normalized vector,
up to one overall phase; taking an absolute square would destroy this
equivalence.

\subsubsection{Three canonical nonchiral examples}
For each example, the displayed polynomial and its normalization
$\mathcal Z$ specify the amplitudes in the common $2^N$-dimensional spin
basis. The production Pfaffian contraction uses this definition without
storing the full vector. The expressions below are at strict infinity.
\begin{enumerate}
\item \textbf{Vacuum $\ket{00}$.}
Choose $(a,b)=(0,0)$, or momentum/winding $(m,w)=(0,0)$.
The continuum vertex is the identity, with $(h,\bar h)=(0,0)$.
Neutrality requires $M=N/2$, and
\begin{equation}
 \mathcal P_{0,0}(O;\infty)=
 \prod_{\mu<\nu}(z_{i_\mu}-z_{i_\nu})^2\prod_{\mu=1}^{N/2}z_{i_\mu}.
\end{equation}
This is the half-filled Haldane--Shastry ground-state amplitude.
\item \textbf{Vertex primary $\ket{ss}$.}
Choose $(a,b)=(1,1)$, or $(m,w)=(1,0)$. The continuum field is
$\mathopen{:}\exp[i(\phi_L+\phi_R)/\sqrt2]\mathclose{:}$, with
$(h,\bar h)=(1/4,1/4)$, dimension $1/2$, and conformal spin zero.
Neutrality requires $M=N/2-1$, and
\begin{equation}
 \mathcal P_{1,1}(O;\infty)=
 \prod_{\mu<\nu}(z_{i_\mu}-z_{i_\nu})^2
 \prod_{\mu=1}^{N/2-1}z_{i_\mu}^{2}.
\end{equation}
\item \textbf{Vertex primary $\ket{3s,s}$.}
Choose $(a,b)=(3,1)$, or $(m,w)=(2,1)$. The continuum field is
$\mathopen{:}\exp[i(3\phi_L+\phi_R)/\sqrt2]\mathclose{:}$, with
$(h,\bar h)=(9/4,1/4)$, dimension $5/2$, and conformal spin two.
Neutrality requires $M=N/2-2$, and
\begin{equation}
 \mathcal P_{3,1}(O;\infty)=
 \prod_{\mu<\nu}(z_{i_\mu}-z_{i_\nu})^2
 \prod_{\mu=1}^{N/2-2}z_{i_\mu}^{4}.
\end{equation}
Its status as a Virasoro primary follows from the vertex-field
stress-tensor product in Sec.~\ref{sec:bosonconstruction}. Calling it a descendant of the extended
$SU(2)_1$ current algebra refers to a different symmetry algebra.
\end{enumerate}
For the original data convention, multiply the last two polynomials by
$\prod_{\mu}(1-z_{i_\mu}/10^8)$ before normalization; the vacuum is
unchanged. The three vectors have different occupation numbers and
are exactly orthogonal without a Gram--Schmidt procedure. The total
spin component in the convention $S_j^z=s_j/2$ is
$S^z_{\rm tot}=M-N/2=-(a+b)/2$.

The production data use these same normalized vertex/Jastrow states,
with $W=10^8$ fixed to match the original implementation. For larger
chains, their interval matrices are obtained by the exact Pfaffian
contraction in Sec.~\ref{sec:extension}; this preserves the amplitudes
while avoiding full-chain enumeration. The independent polynomial/block
comparison and the Haldane--Shastry sanity check are given in
Appendix~\ref{sec:boson-sanity}.
